\documentclass{controle}
\usepackage{main}

\title{Évaluation n°9}
\author{Premiere Spécialité Mathématique}
\date{19 Mars 2026}

% \printanswers

\begin{document}
\titleandname{}
\version{}
\begin{questions}
\question Soit $A$ et $B$ deux événements d'une expérience aléatoire d'univers $\Omega$. Rappeler la formule de probabilités composées permettant de calculer $P(A \cap B)$.

\begin{equation*}
P(A \cap B) =
\end{equation*}

\begin{solution}
\begin{equation*}
P(A \cap B) = P(A) P_A(B)
\end{equation*}
\end{solution}

\question Dans une classe de première spécialité (fictive), les élèves arrivent au hasard en cours. On s'interresse à l'identité de la première personne arrivée sur place. On suppose que les élèves arrivent au hasard de façon équiprobable.

\begin{pycode}
from random import *
fille1 = randrange(1,31)
fille2 = randrange(1,31)
fille3 = randrange(1,31)

garcon1 = randrange(1,31)
garcon2 = randrange(1,31)
garcon3 = randrange(1,31)

total_fille = fille1 + fille2 + fille3
total_garcon = garcon1 + garcon2 + garcon3
total1 = fille1 + garcon1
total2 = fille2 + garcon2
total3 = fille3 + garcon3
total = total_fille + total_garcon
\end{pycode}

\begin{center}
\begin{tblr}{hlines, vlines}
\diagbox{Genre}{Classe}&1G1&1G2&1G3&Total\\
Fille&$\py{fille1}$&$\py{fille2}$&$\py{fille3}$&$\py{total_fille}$\\
Garçon&$\py{garcon1}$&$\py{garcon2}$&$\py{garcon3}$&$\py{total_garcon}$\\
Total&$\py{total1}$&$\py{total2}$&$\py{total3}$&$\py{total}$
\end{tblr}
\end{center}


On pose les événements suivants :
\begin{itemize}
\item $F$ : \og Le premier élève arrivé est une fille \fg
\item $U$ : \og Le premier élève arrivé est en 1G1 \fg
\item $D$ : \og Le premier élève arrivé est en 1G2 \fg
\item $T$ : \og Le premier élève arrivé est en 1G3 \fg
\end{itemize}
\begin{parts}
\part Donner les probabilités $P(F)$ et $P(T)$.
\part Donner les probabilités $P_U(\overbar{F})$ et $P_{\overbar{F}}(U)$.
\part Donner les probabilités $P(D \cap F)$ et $P(T \cap \overbar{F})$.
\end{parts}
\begin{solution}
\begin{parts}
\part $P(F) = \dfrac{\py{total_fille}}{\py{total}}$ et $P(T) = \dfrac{\py{total3}}{\py{total}}$.
\part $P_U(\overbar{F}) = \dfrac{\py{garcon1}}{\py{total1}}$ et $P_{\overbar{F}}(U) = \dfrac{\py{garcon1}}{\py{total_garcon}}$.
\part $P(D \cap F) = \dfrac{\py{fille2}}{\py{total}}$ et $P(T \cap \overbar{F}) = \dfrac{\py{garcon3}}{\py{total}}$
\end{parts}
\end{solution}
\end{questions}
\newpage
\titleandname{}
\version{}
\begin{questions}
\question Soit $A$ et $B$ deux événements d'une expérience aléatoire d'univers $\Omega$. Rappeler la formule permettant de calculer $P_A(B)$.

\begin{equation*}
P_A(B) =
\end{equation*}

\begin{solution}
\begin{equation*}
P_A(B) = \dfrac{P(A \cap B)}{P(A)}
\end{equation*}
\end{solution}

\question Dans une classe de première spécialité (fictive), les élèves arrivent au hasard en cours. On s'interresse à l'identité de la première personne arrivée sur place. On suppose que les élèves arrivent au hasard de façon équiprobable.

\begin{pycode}
from random import *
fille1 = randrange(1,31)
fille2 = randrange(1,31)
fille3 = randrange(1,31)

garcon1 = randrange(1,31)
garcon2 = randrange(1,31)
garcon3 = randrange(1,31)

total_fille = fille1 + fille2 + fille3
total_garcon = garcon1 + garcon2 + garcon3
total1 = fille1 + garcon1
total2 = fille2 + garcon2
total3 = fille3 + garcon3
total = total_fille + total_garcon
\end{pycode}

\begin{center}
\begin{tblr}{hlines, vlines}
\diagbox{Genre}{Classe}&1G1&1G2&1G3&Total\\
Fille&$\py{fille1}$&$\py{fille2}$&$\py{fille3}$&$\py{total_fille}$\\
Garçon&$\py{garcon1}$&$\py{garcon2}$&$\py{garcon3}$&$\py{total_garcon}$\\
Total&$\py{total1}$&$\py{total2}$&$\py{total3}$&$\py{total}$
\end{tblr}
\end{center}

On pose les événements suivants :
\begin{itemize}
\item $F$ : \og Le premier élève arrivé est une fille \fg
\item $U$ : \og Le premier élève arrivé est en 1G1 \fg
\item $D$ : \og Le premier élève arrivé est en 1G2 \fg
\item $T$ : \og Le premier élève arrivé est en 1G3 \fg
\end{itemize}
\begin{parts}
\part Donner les probabilités $P(F)$ et $P(T)$.
\part Donner les probabilités $P_U(\overbar{F})$ et $P_{\overbar{F}}(U)$.
\part Donner les probabilités $P(D \cap F)$ et $P(T \cap \overbar{F})$.
\end{parts}
\begin{solution}
\begin{parts}
\part $P(F) = \dfrac{\py{total_fille}}{\py{total}}$ et $P(T) = \dfrac{\py{total3}}{\py{total}}$.
\part $P_U(\overbar{F}) = \dfrac{\py{garcon1}}{\py{total1}}$ et $P_{\overbar{F}}(U) = \dfrac{\py{garcon1}}{\py{total_garcon}}$.
\part $P(D \cap F) = \dfrac{\py{fille2}}{\py{total}}$ et $P(T \cap \overbar{F}) = \dfrac{\py{garcon3}}{\py{total}}$
\end{parts}
\end{solution}
\end{questions}
\end{document}